\chapter{Introdução}
\label{cap:introducao}

% TODO(P1): Estruturar a introdução no modelo em funil clássico:
% 1. Contextualização: importância dos problemas de otimização combinatória na computação e indústria real.
% 2. O Problema / Limitação: ineficiência de algoritmos recursivos ingênuos que sofrem com explosão combinatória
%    devido à recomputação desnecessária de subproblemas sobrepostos.
% 3. Objetivo do Trabalho: apresentar, formalizar e validar empiricamente o paradigma de Programação Dinâmica (Cap. 15 do Cormen).
% 4. Entregáveis do Grupo: relatório técnico-científico, prova de conceito modular em Python/C++, suíte de experimentos e simuladores web.
% 5. Roteiro de Leitura: breve resumo do que cada capítulo deste relatório aborda (Fundamentação, Paradigma, Análise, Estudo de Caso e Conclusões).

% [Texto a ser redigido pelo integrante P1 - ver docs/tarefas/P1_relatorio_secoes_1_2_6.md]
